% Source PDF pages 131--134; manual pages 2-102--2-105.
\subsection{Optimization}\label{sec:optimization}

The general nonlinear programming problem varies $n$ parameters $x_i$ to minimize
an objective function subject to $m$ equality and inequality constraints:
\begin{equation}
\begin{aligned}
  &\underset{x_i}{\operatorname{minimize}}
  && f(x_i), && i=1,2,\ldots,n,\\
  &\text{subject to}
  && c_j(x_i)=0, && j=1,2,\ldots,m_e,\\
  &&& c_j(x_i)\ge 0, && j=m_e+1,\ldots,m.
\end{aligned}
\label{eq:general-nonlinear-program}
\end{equation}
The equality constraints are conventionally listed first, followed by the inequality
constraints.

\taossubhead{Han--Powell Method (vf02)}

TAOS uses the Han--Powell method, originally distributed by Powell as
\texttt{vf02ad}, to solve the nonlinear programming problem.\manualref{3,22} The
method belongs to the recursive quadratic programming family. It solves a sequence
of quadratic subproblems using an active-constraint set. The active set contains all
equality constraints plus selected inequality constraints currently treated as
equalities.

Initial estimates for the parameters are required. At each iteration, the parameters
are modified to reduce the objective function while satisfying the constraints. For a
fixed active set, the nonlinear problem reduces to
\begin{equation}
\begin{aligned}
  &\underset{x}{\operatorname{minimize}} && f(x),\\
  &\text{subject to} && c(x)=0.
\end{aligned}
\label{eq:active-set-equality-problem}
\end{equation}
Each iteration first determines a search direction and then performs a
one-dimensional minimization in that direction.

For a trial increment $\delta x$, the objective function is approximated by
\begin{equation}
  f(x+\delta x)-f(x)
  \simeq Q(\delta x)
  =f_x(x)\,\delta x
   +\frac{1}{2}\delta x^{T}f_{xx}(x)\delta x,
  \label{eq:optimization-quadratic-approximation}
\end{equation}
and the constraints are linearized as
\begin{equation}
  c(x+\delta x)
  \simeq c(x)+c_x(x)\,\delta x.
  \label{eq:optimization-constraint-linearization}
\end{equation}
Here $f_x$ and $c_x$ are gradients and $f_{xx}$ is the Hessian of the objective.
The Hessian is replaced by a positive-definite variable-metric matrix $B$, giving the
quadratic subproblem
\begin{equation}
\begin{aligned}
  &\underset{\delta x}{\operatorname{minimize}}
  && Q(\delta x)
     =f_x(x)\,\delta x
      +\frac{1}{2}\delta x^{T}B\delta x,\\
  &\text{subject to}
  && c_x(x)\,\delta x+c(x)=0.
\end{aligned}
\label{eq:optimization-quadratic-subproblem}
\end{equation}

The solution supplies a search direction $\delta x$. A one-dimensional search then
chooses the step along that direction, using a penalty function to account for
constraint violations. TAOS uses a version of \texttt{vf02ad} obtained from
D.~G. Hull at The University of Texas at Austin and converted to C by the manual's
author.

\taossubhead{Optimization Parameters}

The user designates optimization parameters with the special keywords described in
\cref{sec:block-optimize}. Any numerical problem-file value may be a parameter,
although initial conditions, guidance rules, and segment final conditions are the
most common choices. Optimization is generally more reliable and economical when
the number of parameters is kept small.

Guidance variables control trajectory shape. A common trajectory-shaping problem
uses a tabulated control history, such as angle of attack versus time, with the tabulated
angles and the final time designated as optimization parameters. The intermediate
time coordinates remain fixed, as shown in
\cref{fig:trajectory-shaping-optimization}.

\begin{figure}[htbp]
  \centering
  \begin{tikzpicture}[>=Latex,x=1.0cm,y=0.76cm,line cap=round]
  \draw[->] (0,0)--(9.7,0) node[below] {Time};
  \draw[->] (0,0)--(0,5.15);
  \node[left,align=center] at (-0.18,2.35) {Angle\\of\\attack};
  \draw[line width=0.9pt]
    plot coordinates {(0,1.9) (1.2,3.0) (2.35,3.0) (3.45,2.25) (4.75,1.55) (6.05,2.40) (7.25,1.82)};
  \foreach \x/\lab in {0/{t_0},1.2/{t_1},2.35/{t_2},3.45/{t_3},4.75/{t_4},6.05/{t_5},7.25/{t_x}}
    \node[below] at (\x,0) {$\lab$};
  \draw[densely dashed] (0,0)--(0,1.9);
  \draw[densely dashed] (1.2,0)--(1.2,3.0);
  \draw[densely dashed] (2.35,0)--(2.35,3.0);
  \draw[densely dashed] (3.45,0)--(3.45,2.25);
  \draw[densely dashed] (4.75,0)--(4.75,1.55);
  \draw[densely dashed] (6.05,0)--(6.05,2.40);
  \draw[densely dashed] (7.25,0)--(7.25,1.82);
  \node[align=center,font=\small] (angles) at (3.40,4.55)
    {Angle-of-attack values are\\defined as optimization parameters};
  \draw[->] (angles.south east)--(3.70,2.25);
  \node[align=left,font=\small] (final) at (8.45,3.75)
    {Final time is also\\defined as an\\optimization parameter};
  \draw[->] (final.south west)--(7.30,0.18);
\end{tikzpicture}

  \caption{Optimization used for trajectory shaping.}
  \label{fig:trajectory-shaping-optimization}
\end{figure}

A reasonable preliminary trajectory is important because the optimization is
sensitive to the initial estimates of the final time and control history. The procedure
works best when the control-history points are evenly distributed. Because the final
time is itself variable, it can move too close to the preceding time point or far beyond
the remaining history. TAOS can periodically restart the optimization and
redistribute the time points evenly.

\taossubhead{Objective Functions and Constraints}

The user supplies the objective function and constraints through the problem file.
Common objectives include time, range, and velocity. A maximization is performed
by minimizing the negative of the desired quantity. Equality and inequality
constraints may reference trajectory output values at the end of any segment and may
combine values from multiple trajectories.

Constraints that apply over the entire path require integral user variables, as
described in \cref{sec:trajectory-block-define}. For example, a path-accumulated
altitude violation may be defined by
\begin{equation}
  h_{\max}=
  \begin{cases}
    \displaystyle\int (h-100000)^2\,dt, & h>100000,\\[4pt]
    0, & \text{otherwise},
  \end{cases}
  \label{eq:maximum-altitude-integral}
\end{equation}
so the variable grows whenever altitude exceeds \SI{100000}{\foot}. The source then
states the optimization constraint
\begin{equation}
  \operatorname{Constrain}\quad h_{\max}<0.
  \label{eq:maximum-altitude-constraint}
\end{equation}
The strict inequality is reproduced as printed; in practical use the integral's initial
value and the optimizer's inequality convention determine how the zero-violation
state is represented.

Large differences in the magnitude of objectives and constraints can cause numerical
problems. TAOS therefore permits normalization by reference values chosen so the
scaled quantities are typically between one and ten. Reference values may be
selected automatically or supplied by the user.

\taossubhead{Gradients}

The optimizer requires the gradients $f_x(x)$ and $c_x(x)$. TAOS computes them
numerically with forward or central differences. For the same perturbation size,
central differences are normally more accurate but require an additional function
evaluation. Forward differences can often obtain adequate accuracy with a smaller
perturbation and are less expensive; both choices are available.
